\chapter*{Introdução}

Sistemas Especialistas (SEs) são sistemas computacionais de inteligência artificial concebidos para resolver problemas em um domínio específico, utilizando uma representação explícita do conhecimento especializado e mecanismos de inferência para produzir conclusões, diagnósticos ou recomendações semelhantes aos que um especialista humano produziria \cite{giarratano1998expert}.

Grandes Modelos de Linguagem, ou, em inglês, \textit{Large Language Models} (LLMs), são modelos computacionais baseados, em geral, em redes neurais profundas, pré-treinados em grandes volumes de dados textuais para aprender padrões estatísticos e representações da linguagem, permitindo realizar tarefas como compreender e gerar textos, responder a perguntas, resumir documentos e produzir código \cite{minaee2024large}.

Associar o uso de LLMs com o Desenvolvimento de Sistemas Especialistas pode trazer  algumas vantagens \cite{yang2025comprehensive, garridomerchan2026gofai}:

\begin{itemize}
    \item Melhor contextualização: fornecer ao LLM conhecimento específico do domínio em que deve atuar;
    \item Maior precisão factual: apoiar as respostas do modelo em fontes de conhecimento estruturadas através, por exemplo, de \textit{prompts} especializados;
    \item Melhor aproveitamento dos recursos de conhecimento: permitir que informações explícitas sejam utilizadas na geração das respostas;
    \item Melhor raciocínio simbólico: o motor de inferência, parte importante de SEs, pode executar deduções sobre a base de conhecimento.
\end{itemize}

Se faz necessário então cobrir de forma prática esta ponte e este livro é dedicado a isso. Ele traz exemplos didáticos de união entre grandes modelos, \text{i.e.} modelos da Google AI, e SEs desenvolvidos pelos discentes da disciplina, Sistemas Especialistas para a Web, turma de 2026.01 do Programa de Pós-Graduação em Desenvolvimento Web do Instituto Federal da Bahia (IFBA), campus Vitória da Conquista.

Nos capítulos seguintes, além de descrições sobre as características de cada solução-exemplo, também são disponibilizados, de forma pública e sob licença permissiva, códigos-fonte e vídeos de apresentação. O intuito é fornecer um material didático para quem, porventura, deseje entender como funcionam SEs e como criá-los utilizando os recursos modernos associados à componentização de software sobre LLMs: \textit{embeddings}, banco de dados vetorial, cache semântico, julgamento automatizado por IA \cite{azevedo2026julgamento} e integração com ferramentas MCP (\textit{Model Context Protocol}).

\begin{flushright}
Luis Paulo da Silva Carvalho.
\end{flushright}